% GENERATED FILE. Edit canonical lessons, then run npm run generate:books.
% canonical-source-hash: fnv1a64:218408c460f62ce2
% canonical-lessons: GU-C72-big, GU-C72-small, GU-C72-long, GU-C72-short, GU-C72-heavy

\chapter{Big and Small, Long and Short}
\label{ch:gu-big-and-small-long-and-short}
\hlchaptermodality{72}

\begin{chapteropening}
\textbf{By the end of this chapter:} \emph{I can say big, small, long, short and heavy.}

Complete the last lesson of chapter 72: I can say big, small, long, short and heavy.
\end{chapteropening}

\section[\texorpdfstring{\gu{મોટું} (moṭũ)}{મોટું (moṭũ)}]{\gu{મોટું} (moṭũ) — big}

\label{lesson:GU-C72-big}



\begin{quote}
Before the new one: say the Gujarati for beautiful, then the Gujarati for bad.
\end{quote}

\subsection*{You'll want to know: \gu{મોટું}}
\textbf{\gu{મોટું}} (\emph{moṭũ}) — \textquotedblleft{}big\textquotedblright{}.

\begin{cousinweb}
CDIAL lists \textbf{\gu{મોટું}} (\emph{moṭũ}) under \emph{*mōṭṭa}, one member of a large group of reconstructed forms it heads \emph{*muṭṭa}, \textquotedblleft{}defective\textquotedblright{} (CDIAL 10187).
\end{cousinweb}

\begin{grammarlens}[title={What you've built}]
One more word for big and small, long and short.
\end{grammarlens}

\subsection*{Your turn}
\begin{itemize}
\raggedright
  \item \emph{Say it:} \emph{moṭũ}
  \item \emph{Say it:} \emph{moṭũ}, once more
  \item \emph{Say it:} say \emph{moṭũ}
  \item \emph{Recall:} say the Gujarati for beautiful, then the Gujarati for bad, then say \emph{moṭũ} again
\end{itemize}

\begin{tcolorbox}[breakable,colback=teal!4,colframe=teal!35!black,title={Before you move on}]
What does \gu{મોટું} mean? (\textquotedblleft{}Big\textquotedblright{}.) Say it once more.
\end{tcolorbox}
\section[\texorpdfstring{\gu{નાનું} (nānũ)}{નાનું (nānũ)}]{\gu{નાનું} (nānũ) — small}

\label{lesson:GU-C72-small}



\begin{quote}
Before the new one: say the Gujarati for bad, then the Gujarati for big.
\end{quote}

\subsection*{You'll want to know: \gu{નાનું}}
\textbf{\gu{નાનું}} (\emph{nānũ}) — \textquotedblleft{}small\textquotedblright{}.

\begin{cousinweb}
\textbf{\gu{નાનું}} (\emph{nānũ}) is a tadbhava. CDIAL lists it under Sanskrit \emph{ślakṣṇa}, \textquotedblleft{}slippery, smooth, soft, tender\textquotedblright{}, which the old word-lists also gloss \textquotedblleft{}thin, small\textquotedblright{} (CDIAL 12732).
\end{cousinweb}

\begin{grammarlens}[title={What you've built}]
One more word for big and small, long and short.
\end{grammarlens}

\subsection*{Your turn}
\begin{itemize}
\raggedright
  \item \emph{Say it:} \emph{nānũ}
  \item \emph{Say it:} \emph{nānũ}, once more
  \item \emph{Say it:} say \emph{nānũ}
  \item \emph{Recall:} say the Gujarati for bad, then the Gujarati for big, then say \emph{nānũ} again
\end{itemize}

\begin{tcolorbox}[breakable,colback=teal!4,colframe=teal!35!black,title={Before you move on}]
What does \gu{નાનું} mean? (\textquotedblleft{}Small\textquotedblright{}.) Say it once more.
\end{tcolorbox}
\section[\texorpdfstring{\gu{લાંબું} (lā̃bũ)}{લાંબું (lā̃bũ)}]{\gu{લાંબું} (lā̃bũ) — long, tall}

\label{lesson:GU-C72-long}



\begin{quote}
Before the new one: say the Gujarati for big, then the Gujarati for small.
\end{quote}

\subsection*{You'll want to know: \gu{લાંબું}}
\textbf{\gu{લાંબું}} (\emph{lā̃bũ}) — \textquotedblleft{}long, tall\textquotedblright{}.

\begin{cousinweb}
\textbf{\gu{લાંબું}} (\emph{lāṁbũ}) is a tadbhava of Sanskrit \emph{lamba}, \textquotedblleft{}pendent\textquotedblright{}, which already meant \textquotedblleft{}long\textquotedblright{} in compounds (CDIAL 10951).
\end{cousinweb}

\begin{grammarlens}[title={What you've built}]
One more word for big and small, long and short.
\end{grammarlens}

\subsection*{Your turn}
\begin{itemize}
\raggedright
  \item \emph{Say it:} \emph{lā̃bũ}
  \item \emph{Say it:} \emph{lā̃bũ}, once more
  \item \emph{Say it:} say \emph{lā̃bũ}
  \item \emph{Recall:} say the Gujarati for big, then the Gujarati for small, then say \emph{lā̃bũ} again
\end{itemize}

\begin{tcolorbox}[breakable,colback=teal!4,colframe=teal!35!black,title={Before you move on}]
What does \gu{લાંબું} mean? (\textquotedblleft{}Long, tall\textquotedblright{}.) Say it once more.
\end{tcolorbox}
\section[\texorpdfstring{\gu{ટૂંકું} (ṭū̃kũ)}{ટૂંકું (ṭū̃kũ)}]{\gu{ટૂંકું} (ṭū̃kũ) — short}

\label{lesson:GU-C72-short}



\begin{quote}
Before the new one: say the Gujarati for small, then the Gujarati for long.
\end{quote}

\subsection*{You'll want to know: \gu{ટૂંકું}}
\textbf{\gu{ટૂંકું}} (\emph{ṭū̃kũ}) — \textquotedblleft{}short\textquotedblright{}.

\begin{cousinweb}
CDIAL lists \textbf{\gu{ટૂંકું}} (\emph{ṭūṁkũ}) under \emph{*ṭukka}, \textquotedblleft{}piece\textquotedblright{}, a form it reconstructs beside \emph{*ṭukk}, \textquotedblleft{}cut, break\textquotedblright{} (CDIAL 5466). It thinks a link with Sanskrit \emph{truṭ}, \textquotedblleft{}to break\textquotedblright{}, very doubtful.
\end{cousinweb}

\begin{grammarlens}[title={What you've built}]
One more word for big and small, long and short.
\end{grammarlens}

\subsection*{Your turn}
\begin{itemize}
\raggedright
  \item \emph{Say it:} \emph{ṭū̃kũ}
  \item \emph{Say it:} \emph{ṭū̃kũ}, once more
  \item \emph{Say it:} say \emph{ṭū̃kũ}
  \item \emph{Recall:} say the Gujarati for small, then the Gujarati for long, then say \emph{ṭū̃kũ} again
\end{itemize}

\begin{tcolorbox}[breakable,colback=teal!4,colframe=teal!35!black,title={Before you move on}]
What does \gu{ટૂંકું} mean? (\textquotedblleft{}Short\textquotedblright{}.) Say it once more.
\end{tcolorbox}
\section[\texorpdfstring{\gu{ભારે} (bhāre)}{ભારે (bhāre)}]{\gu{ભારે} (bhāre) — heavy}

\label{lesson:GU-C72-heavy}



\begin{quote}
Before the new one: say the Gujarati for long, then the Gujarati for short.
\end{quote}

\subsection*{You'll want to know: \gu{ભારે}}
\textbf{\gu{ભારે}} (\emph{bhāre}) — \textquotedblleft{}heavy\textquotedblright{}.

\begin{cousinweb}
\textbf{\gu{ભારે}} (\emph{bhāre}) is built on Gujarati \emph{bhār}, \textquotedblleft{}load, weight\textquotedblright{}, which CDIAL lists under Sanskrit \emph{bhāra}, \textquotedblleft{}burden, weight\textquotedblright{} (CDIAL 9459). It lists Gujarati \emph{bhārī}, \textquotedblleft{}heavy\textquotedblright{}, under \emph{bhārika}, \textquotedblleft{}forming a load, heavy\textquotedblright{} (CDIAL 9465).
\end{cousinweb}

\begin{grammarlens}[title={What you've built}]
One more word for big and small, long and short.
\end{grammarlens}

\subsection*{Your turn}
\begin{itemize}
\raggedright
  \item \emph{Say it:} \emph{bhāre}
  \item \emph{Say it:} \emph{bhāre}, once more
  \item \emph{Say it:} say \emph{bhāre}
  \item \emph{Recall:} say the Gujarati for long, then the Gujarati for short, then say \emph{bhāre} again
\end{itemize}

\begin{tcolorbox}[breakable,colback=teal!4,colframe=teal!35!black,title={Before you move on}]
What does \gu{ભારે} mean? (\textquotedblleft{}Heavy\textquotedblright{}.) Say it once more.
\end{tcolorbox}
